\documentclass[12pt]{article}
\usepackage[utf8]{inputenc}
\usepackage[T1]{fontenc}
\usepackage{hyperref}
\title{Z\'aparo Language Resources for K-12 Curriculum Development}
\author{Compiled by Claude Sonnet 5.5 (claude-sonnet-5-5) for Andr\'as Kornai}
\date{}
\begin{document}
\maketitle

\section{Introduction}
Z\'aparo (ISO 639-3 zro; Glottocode zapa1253) is a Zaparoan language of Ecuador and Peru. It is nearly extinct. Sources report only a handful of elderly speakers, and most of the remaining ethnic Z\'aparo speak Quechua. Revitalization work has been under way since 1999. The Z\'aparo organization in Ecuador prepared a linguistic description for developing teaching materials. It has used these in community primary schools and in an adult literacy programme. A UNESCO-funded project ran in 2005--2007. Some details in the sources are dated and inconsistent.

\section{Primary Educational Resources}
No standalone published textbook or primer was identified in the search results. The revitalization programme is reported to have produced teaching materials, but their text was not located online.

\section{Linguistic and Reference Materials}
A document by Viatori on Z\'aparo is held in the AILLA archive: \url{https://ailla-legacy.lib.utexas.edu/sites/default/files/documents/Viatori_Zapara.pdf}. Omniglot has an overview of the language and its writing: \url{https://mail.omniglot.com/writing/zaparo.htm}. A Wikipedia overview with linguistic and speaker information is at \url{https://wikipedia.com/wiki/ISO_639:zro}.

\section{Religious and Community Resources}
UNESCO documents a bi-national (Ecuador and Peru) action plan for safeguarding and revitalizing the oral and intangible heritage of the Z\'aparo community: \url{https://ich.unesco.org/en/projects/bi-national-action-plan-for-the-safeguarding-and-revitalization-of-the-oral-and-intangible-heritage-of-the-zapara-community-00012}. The book \emph{One State, Many Nations} appeared in the results, possibly as context on Indigenous nations in Ecuador: \url{https://ubcpress.ca/one-state-many-nations}. Its relevance to Z\'aparo was not verified.

\section{Digital and Audio Resources}
No dedicated audio or video collections were found in the searches. The AILLA archive may hold further material.

\section{Curriculum Development Context}
Because there are very few speakers, a curriculum would be a revitalization and heritage curriculum built on existing documentation, the community's earlier school materials and the archived recordings. Any work should involve the Z\'aparo organization and the elders. Current speaker numbers vary between sources.
\end{document}